\chapter{Systemarchitektur}
\label{ch:architektur}

% ==============================================================================
\section{Ausgangslage und Idee}
\label{sec:arch-idee}

Als die Arbeit an Modul~4 begann, standen weder der Roboter noch die Kameras zur Verfügung.
Gleichzeitig ist der teuerste Fehler bei Behavior Cloning einer, der erst nach der
Datenaufnahme auffällt: Er entwertet den ganzen Datensatz (\cref{sec:unumkehrbar}). Daraus
ergaben sich zwei Anforderungen mit derselben Lösung (N2 und N3 aus
\cref{sec:anf-nichtfunktional}): Die gesamte Logik soll ohne Hardware laufen und getestet
sein, und Roboter und Kameras sollen austauschbar sein. Umgesetzt wurde das mit einer
Architektur aus Ports und Adaptern.

% ==============================================================================
\section{Schichtenmodell mit Ports und Adaptern}
\label{sec:ports}

Die Software ist in vier Schichten gegliedert (\cref{fig:schichten}). Die
\emph{Pipeline} enthält alle Logik von der Bahnplanung über das Rauschen bis zur Ausführung
der Policy. Sie spricht Hardware ausschließlich über abstrakte Schnittstellen an, die
\emph{Ports}: \code{RobotPort} für den Roboter, \code{CameraPort} für Kameras,
\code{TeleopPort} für Eingabegeräte beim Teachen und \code{ClockPort} für die Zeit. Für jeden
Port gibt es mehrere \emph{Adapter}: einen für die reale Hardware und einen oder mehrere für
die Simulation.

\begin{figure}[htbp]
  \centering
  \begin{tikzpicture}[font=\small,
      lay/.style={draw, rounded corners=1.5mm, minimum height=10mm, text width=108mm,
                  align=left, inner xsep=3mm},
      tag/.style={font=\small\bfseries, anchor=east, text width=22mm, align=right}]
    \node[lay, fill=black!5] (app) at (0,0)
      {\code{teach}\quad \code{record}\quad \code{export}\quad \code{train}\quad \code{infer}\quad
       \code{eval}\quad Bedienoberfläche};
    \node[lay, fill=blue!8, minimum height=16mm, below=3mm of app] (pipe)
      {Bahnplanung, Rauschen, Kollisionsprüfung, Synchronisation, Aufnahme, Datensatz,
       Policy-Ausführung, Sicherheitsprüfung, Auswertung\\[1mm]
       \scriptsize importiert nie \code{neurapy} oder \code{gxipy} (durch Test erzwungen)};
    \node[lay, fill=orange!12, below=3mm of pipe] (ports)
      {\code{RobotPort}\quad \code{CameraPort}\quad \code{TeleopPort}\quad \code{ClockPort}};
    \node[lay, fill=green!8, minimum height=16mm, below=3mm of ports] (ad)
      {real:\quad Neura (NeuraPy), Daheng (\code{gxipy}), UVC-Kamera, Gamepad, Echtzeit\\
       simuliert:\quad SimRobot (\acs{urdf}), Simulations-, Replay- und Dateikameras,
       Skript-Teleop, Simulationsuhr};
    \node[tag] at ([xshift=-2mm]app.west) {Anwendungen};
    \node[tag] at ([xshift=-2mm]pipe.west) {Pipeline};
    \node[tag] at ([xshift=-2mm]ports.west) {Ports};
    \node[tag] at ([xshift=-2mm]ad.west) {Adapter};
  \end{tikzpicture}
  \caption[Schichtenmodell der Software]{Schichtenmodell der Software. Die Pipeline kennt
  nur die Ports; welcher Adapter dahinter steht, wird beim Start festgelegt.}
  \label{fig:schichten}
\end{figure}

Die Trennung ist nicht nur eine Konvention. Ein automatischer Test durchsucht alle
Pipeline-Module und schlägt fehl, sobald eines davon \code{neurapy} oder \code{gxipy}
importiert, auch verzögert innerhalb einer Funktion. Zeit und Zufall sind ebenfalls über
Ports eingebunden: Mit einer simulierten Uhr und festen Startwerten der Zufallsgeneratoren
laufen Tests reproduzierbar und schneller als in Echtzeit.

\paragraph{Was die Trennung leistet und was nicht.}
Abgesichert wird damit die \emph{Logik}: Synchronisationsregeln, Zuordnung von Bild und
Aktion, Bahn- und Rauschmathematik, Kollisionsprüfung, Datensatzformat, Trainings- und
Ausführungsschleife sowie alle Fehler- und Abbruchpfade. Nicht abgesichert werden die
\emph{Physik} und die tatsächlichen Eigenschaften der Schnittstellen: reale Latenzen,
\acs{usb}-Bandbreite, das Verhalten der Steuerung bei \code{servo_j}, die Totzeit des
Greifers und die Bildstatistik. Ein simulierter Adapter bildet immer nur die eigene Annahme
über die Hardware ab. Das Risiko verschwindet also nicht, es verlagert sich in die Adapter.
Deshalb wurden die Adapter an der Hardware gesondert abgenommen (\cref{sec:testkonzept}).

% ==============================================================================
\section{Modulübersicht}
\label{sec:module}

\Cref{tab:module} zeigt die wichtigsten Module. Die Anwendungen in \code{apps/} sind
Kommandozeilenprogramme, die auch die Bedienoberfläche aufruft. \code{tools/} enthält
Prüf- und Auswertewerkzeuge, die nicht Teil des Ablaufs sind.

\begin{table}[htbp]
  \centering
  \caption{Wichtige Module der Pipeline (\code{bc/})}
  \label{tab:module}
  \small
  \begin{tabularx}{\textwidth}{@{}l X@{}}
    \toprule
    Modul & Aufgabe \\
    \midrule
    \code{ports.py}, \code{adapters/} & Schnittstellen und ihre realen und simulierten
      Umsetzungen \\
    \code{config.py} & alle Konstanten an einer Stelle, mit Begründung im Kommentar \\
    \code{sequence.py}, \code{trajectory.py} & Ablauf als Punktfolge, Bahnplanung mit
      Rampen und Überschleifen \\
    \code{noise.py} & Rauschprozess, Trichter, Versatz im Anflug, Neuziehen ungültiger
      Bahnen \\
    \code{collision.py}, \code{table.py}, \code{obstacles.py} & Kollisionsmodell aus Tisch,
      Sperrzonen und festen Hindernissen \\
    \code{sync.py}, \code{capture.py}, \code{recorder.py} & Bild- und Zustandsabgriff,
      Synchronisation, Aufzeichnung einer Episode \\
    \code{put_back.py}, \code{auto_layout.py} & Zurücklegen und automatisches Umlegen des
      Bauteils \\
    \code{dataset.py}, \code{lerobot_io.py} & eigenes Aufnahmeformat und Export nach
      LeRobot \\
    \code{diffusion.py}, \code{policy.py} & Laden der Diffusion Policy, Ausführung von
      Aktionsfolgen \\
    \code{execution.py}, \code{servo.py}, \code{live.py} & Taktschleife und
      Sollwertstrom zum Roboter \\
    \code{safety.py} & Prüfung jedes Sollwerts, Arbeitsraumüberwachung, Stopp \\
    \code{gui/} & Reiter von Modul~4 in der Bedienoberfläche \\
    \bottomrule
  \end{tabularx}
\end{table}

% ==============================================================================
\section{Kinematik}
\label{sec:kinematik}

Bahnplanung, Rauschen und Kollisionsprüfung brauchen inverse und Vorwärtskinematik. Bei
der LARA~5 liegen beide in der Steuerung und sind nur über NeuraPy erreichbar. Damit wäre
genau der aufwändigste und fehleranfälligste Teil, die Erzeugung verrauschter Bahnen, ohne
Roboter nicht ausführbar gewesen.

\paragraph{Idee.} Aus dem \ac{urdf}-Modell, das mit der Steuerungssoftware ausgeliefert
wird, wurde eine eigene Kinematik für den simulierten Roboter gebaut: Vorwärtskinematik
analytisch entlang der Gelenkkette, inverse Kinematik numerisch mit gedämpfter
Pseudoinverse,
\begin{equation}
  \Delta q = J^{\mathsf T}\bigl(J J^{\mathsf T} + \lambda^2 I\bigr)^{-1} e ,
  \label{eq:dls}
\end{equation}
wobei $J$ die Jacobi-Matrix, $e$ der Posenfehler und $\lambda$ ein Dämpfungsfaktor ist, der
die Lösung in der Nähe singulärer Stellungen stabil hält.

\paragraph{Problem.} Eine Gegenprobe mit einer real gemessenen Stellung traf die Pose der
Steuerung mit keinem der ausgelieferten Modelle; der kleinste Fehler lag bei etwa
\SI{83}{\milli\metre}. Das \ac{urdf} ist damit als Strukturvorlage brauchbar, aber nicht als
verbindliches Modell der Anlage.

\paragraph{Lösung.} Es wurde zweigleisig gearbeitet. Die eigene Kinematik dient nur dazu,
die \emph{Logik} offline zu prüfen; Zahlen wie Erreichbarkeit oder Abstände aus der
Simulation werden nicht als Nachweis verwendet. An der virtuellen Steuerung und später an der
Anlage rechnet die Pipeline ausschließlich mit der Kinematik der Steuerung. Ein
Prüfwerkzeug vergleicht dazu Vorwärtskinematik, inverse Kinematik und die gemeldete
\ac{tcp}-Pose. An der Anlage stimmten Vorwärtskinematik und \ac{vm} exakt überein, und die
Rückprobe der inversen Kinematik ergab keinen Fehler \quelle{Berichte/2026-10-05_12-59_Kinematik-Anlage.json}.
Damit ist die \ac{vm} für Bahnplanung und Kollisionsprüfung gleichwertig zur Anlage.

Bei der inversen Kinematik wird in jedem Schritt die Lösung des vorherigen Schritts als
Startwert übergeben. Die Steuerung liefert dann die nächstgelegene Lösung, und Sprünge
zwischen Armkonfigurationen sind ausgeschlossen. Ein Vortest an drei realen Posen ergab, dass
alle Punkte in einem Korridor von \SI{\pm 20}{\milli\metre} erreichbar sind, ohne
Konfigurationssprung und weit von singulären Stellungen entfernt, bei
\SI{2.2}{\milli\second} je Aufruf.

% ==============================================================================
\section{Testkonzept}
\label{sec:testkonzept}

Die Tests bauen in drei Stufen aufeinander auf:
\begin{enumerate}
  \item \textbf{Unit-Tests} prüfen einzelne Module hardwarefrei und deterministisch, zum
    Beispiel Rauschprozess, Bahnplanung, Kollisionsmodell und Datensatzformat.
  \item \textbf{Contract-Tests} prüfen die Zusagen eines Ports, also Wertebereiche,
    Einheiten, Bezugssysteme, Zeitstempel und Verhalten im Fehlerfall. Dieselbe Testsuite
    läuft gegen den simulierten und gegen den realen Adapter. Aus \enquote{der Adapter ist
    noch ungetestet} wird so eine abhakbare Liste für den ersten Tag an der Anlage.
  \item \textbf{End-to-End-Läufe} in der Simulation führen die ganze Kette aus: Punkte,
    Bahnplanung mit Rauschen, Aufnahme, Export, Kurztraining und Ausführung der Policy. Sie
    weisen nach, dass Zustand und Aktion bei Aufnahme und Ausführung zusammenpassen.
\end{enumerate}

Damit die simulierten Adapter nicht nur zeigen, dass nichts abstürzt, lassen sie sich gezielt
fehlerhaft einstellen: Latenz und Schwankung, verlorene und veraltete Bilder, fehlende
Lösungen der inversen Kinematik, leere Antworten der Steuerung, ausgelöster Schutzstopp und
verschobene Kamera. Erst damit werden die Abbruch- und Verwerfungspfade tatsächlich
durchlaufen. Die Testsuite umfasst derzeit rund 600 Testfälle und läuft auch ohne
\code{pytest} (\code{tests/run_all.py}).

Für den ersten Zugang zur Anlage wurde aus dem, was hardwarefrei nicht prüfbar ist, eine
Abnahmeliste abgeleitet: erreichte Rate und Zeitversatz je Quelle, Formate und Einheiten der
NeuraPy-Aufrufe, Roundtrip von \code{servo_j}, Kinematik, Greifer-Totzeit,
Werkzeuggeometrie, beide Kameras, Kollisionsmodell und Not-Halt. Die zugehörigen
Prüfwerkzeuge (\code{tools/check_*.py}) schreiben ihre Ergebnisse als Bericht in den
Ordner \code{Berichte/}.

% ==============================================================================
\section{Simulation und virtuelle Steuerung}
\label{sec:sim-vm}

Modul~4 wurde in drei Umgebungen erprobt, die jeweils einen anderen Teil absichern
(\cref{tab:umgebungen}).

\begin{table}[htbp]
  \centering
  \caption{Umgebungen und was sie absichern}
  \label{tab:umgebungen}
  \small
  \begin{tabularx}{\textwidth}{@{}l X X@{}}
    \toprule
    Umgebung & sichert ab & sichert nicht ab \\
    \midrule
    Eigene Simulation (\code{SimRobot}) & gesamte Logik, schneller als Echtzeit, mit
      Fehlerinjektion & Geometrie der Anlage, Bilder, Verhalten der Steuerung \\
    \acs{vm} der Neura-Steuerung & NeuraPy-Schnittstelle, Kinematik, Sollwertstrom,
      Bahnplanung mit echten Punkten & Bilder, Greifer, Mechanik (der Arm folgt ideal) \\
    Reale Anlage & alles, insbesondere Kameras, Schwingungen und Greifen & -- \\
    \bottomrule
  \end{tabularx}
\end{table}

Die \ac{vm} wird von Neura als virtuelle Maschine bereitgestellt und spricht dieselbe
Schnittstelle wie die reale Steuerung. Sie läuft auf dem Betriebsrechner mit einem
eigenen Host-only-Netzwerk, weil NeuraPy eine feste Adresse der Steuerung erwartet. Für
Bildtests ohne Anlage gibt es zusätzlich eine Wiedergabe-Kamera, die aufgezeichnete Bilder
einer Episode synchron zur Bahn abspielt. Damit lässt sich eine Policy in der \ac{vm} mit
echten Bildern prüfen (\cref{ch:versuche}).
